\section{均匀连续与Lipschitz连续}
\label{sec:mmap-uniform-lipschitz}

连续性允许$\delta$随输入点改变。处理单个点附近的问题时这无关紧要，但处理柯西列、稠密集上的延拓或迭代误差时，点列的各项散布在不同位置，需要一个对整个定义域都有效的阈值，最好还能给出输出误差是输入误差的多少倍。本节依次加强条件：均匀连续要求阈值与位置无关，Lipschitz连续要求误差按固定比例放大。最后说明紧定义域上连续自动升级为均匀连续，却不能自动升级为Lipschitz连续。卷一在实轴区间上讨论过这些概念，可参考\BookChapterRef{chap:limits-continuity}{卷一“极限与连续”章}，这里推广到度量空间。

\subsection{均匀连续}
\label{sec:mmap-uniform-continuity}

先看阈值怎样依赖位置。在$\mathbb R$上取$f(x)=x^2$，有$\lvert f(x')-f(x)\rvert=\lvert x'-x\rvert\,\lvert x'+x\rvert$。在$x$附近要使输出误差小于$1$，输入误差大约要小于$1/(2\lvert x\rvert)$，$x$越大，允许的输入误差越小。换到区间$[0,1]$上，$\lvert x'+x\rvert\leq2$，于是$\lvert f(x')-f(x)\rvert\leq2\lvert x'-x\rvert$，取$\delta=\varepsilon/2$对区间中所有点都有效。

\begin{definition}[均匀连续]
\label{def:mmap-uniform-continuity}
若对任意$\varepsilon>0$，存在$\delta>0$，使得对所有满足$d_X(x,x')<\delta$的$x,x'\in X$，都有$d_Y(f(x),f(x'))<\varepsilon$，则称$f$\term{均匀连续}（uniformly continuous）。卷一称之为一致连续。
\end{definition}

与逐点连续相比，变化只在量词的次序：逐点连续是“对每个$x$和$\varepsilon$，存在$\delta$”，均匀连续是“对每个$\varepsilon$，存在$\delta$，对所有$x,x'$”。否定均匀连续时，常用点对序列：若存在$\varepsilon>0$及点列$x_n,x_n'$，使$d(x_n,x_n')\to0$而$d(f(x_n),f(x_n'))\geq\varepsilon$，则$f$不均匀连续，因为任何$\delta$都会被某个$d(x_n,x_n')<\delta$的点对打破。对$\mathbb R$上的$x^2$，取$x_n=n$，$x_n'=n+1/n$，则$\lvert x_n-x_n'\rvert=1/n\to0$，而$\lvert x_n'^2-x_n^2\rvert=2+1/n^2\geq2$。

均匀连续的直接用处是保持柯西列。

\begin{theorem}[均匀连续映射保持柯西列]
\label{thm:mmap-uniform-cauchy-image}
若$f$均匀连续，$(x_n)$是$X$中的柯西列，则$(f(x_n))$是$Y$中的柯西列。
\end{theorem}

\begin{proof}
对$\varepsilon>0$取均匀连续的$\delta$，再取$N$使$m,n\geq N$时$d(x_m,x_n)<\delta$，则这些$m,n$满足$d(f(x_m),f(x_n))<\varepsilon$。
\end{proof}

仅有连续性不够。$f(x)=1/x$在$(0,1)$上连续，$x_n=1/(n+1)$是柯西列，像$n+1$却不是柯西列；按定理，这也说明$1/x$在$(0,1)$上不均匀连续。第~\ref{sec:mmap-extension}节把均匀连续映射从稠密子集延拓出去，依靠的正是这一定理：逼近同一点的柯西列，其像也是柯西列。

\subsection{Lipschitz连续}
\label{sec:mmap-lipschitz}

均匀连续只说存在$\delta$，没有说$\delta$与$\varepsilon$的关系。若输出误差不超过输入误差的固定倍数，误差估计就可以直接写出来。

\begin{definition}[Lipschitz连续]
\label{def:mmap-lipschitz}
若存在有限常数$L\geq0$，使对所有$x,x'\in X$有
\begin{equation}
d_Y(f(x),f(x'))\leq L\,d_X(x,x'),
\label{eq:mmap-lipschitz}
\end{equation}
则称$f$为\term{Lipschitz连续}（Lipschitz continuous）的，$L$称为一个Lipschitz常数。满足式\eqref{eq:mmap-lipschitz}的最小常数记为
\begin{equation*}
\operatorname{Lip}(f)=\sup_{x\neq x'}\frac{d_Y(f(x),f(x'))}{d_X(x,x')},
\end{equation*}
称为$f$的\term{Lipschitz常数}（Lipschitz constant）。
\end{definition}

Lipschitz连续蕴含均匀连续：$L>0$时取$\delta=\varepsilon/L$，$L=0$时$f$为常值。它还给出可以计算的误差放大界：输入误差为$\eta$时，输出误差至多为$L\eta$。若$f:X\to Y$与$g:Y\to Z$都Lipschitz连续，则
\begin{equation*}
d(g(f(x)),g(f(x')))\leq\operatorname{Lip}(g)\,d(f(x),f(x'))\leq\operatorname{Lip}(g)\operatorname{Lip}(f)\,d(x,x'),
\end{equation*}
所以$\operatorname{Lip}(g\circ f)\leq\operatorname{Lip}(g)\operatorname{Lip}(f)$：多步计算的误差放大倍数至多相乘。

常见的Lipschitz映射有几类。距离函数$x\mapsto d(x,A)$的Lipschitz常数不超过$1$。区间上可微且$\lvert f'\rvert\leq L$的实函数满足$\lvert f(x)-f(x')\rvert=\lvert f'(\xi)\rvert\,\lvert x-x'\rvert\leq L\lvert x-x'\rvert$，这里用到中值定理，可参考\BookChapterRef{chap:univariate-differentiation}{卷一“一元函数的微分”章}；前面的$x^2$在$[0,1]$上就有$L=2$。矩阵$\bm A$定义的线性映射满足$\lVert\bm A\bm x-\bm A\bm y\rVert\leq\lVert\bm A\rVert\,\lVert\bm x-\bm y\rVert$，其中$\lVert\bm A\rVert$是诱导的算子范数，可参考\BookChapterRef{chap:normed-inner-product-spaces}{卷二“赋范空间与内积空间”章}。

均匀连续不必Lipschitz连续。在$[0,1]$上取$f(x)=\sqrt x$。对$0\leq y\leq x$，$(\sqrt y+\sqrt{x-y})^2=x+2\sqrt{y(x-y)}\geq x$，所以$\sqrt x-\sqrt y\leq\sqrt{x-y}$，取$\delta=\varepsilon^2$即得均匀连续。但$\sqrt x/x=1/\sqrt x$在$x\to0^+$时无界，没有有限的$L$满足$\sqrt x-\sqrt0\leq Lx$。平方根在$0$附近把输入误差$\eta$放大为$\sqrt\eta$，放大倍数$1/\sqrt\eta$没有上界。

$L<1$的自映射把任意两点的距离按固定比例缩小，这是第~\ref{sec:mmap-contraction}节压缩迭代的出发点。

\subsection{紧空间上的均匀连续性}
\label{sec:mmap-compact-uniform}

前面的反例$x^2$与$1/x$分别定义在无界的$\mathbb R$和不包含端点的$(0,1)$上，都不是紧空间。紧性恰好排除了阈值随位置无限缩小的情形。

\begin{theorem}[紧空间上的均匀连续性]
\label{thm:mmap-compact-uniform}
若$X$是紧致度量空间，$f:X\to Y$连续，则$f$均匀连续。
\end{theorem}

\begin{proof}
若不然，存在$\varepsilon>0$及点列$x_n,x_n'$，使$d(x_n,x_n')<1/n$而$d(f(x_n),f(x_n'))\geq\varepsilon$。按第~\ref{chap:metric-spaces}章，紧致度量空间序列紧致，取子列$x_{n_k}\to x\in X$，则$d(x'_{n_k},x)\leq d(x'_{n_k},x_{n_k})+d(x_{n_k},x)\to0$。由$f$在$x$处连续及序列判据，$f(x_{n_k})$与$f(x'_{n_k})$都趋于$f(x)$，于是$d(f(x_{n_k}),f(x'_{n_k}))\to0$，与$\geq\varepsilon$矛盾。
\end{proof}

也可以用开覆盖证明：对每个$x$取使$f(B(x,\delta_x))\subseteq B(f(x),\varepsilon/2)$的$\delta_x$，球族$\{B(x,\delta_x)\}$的Lebesgue数就是对全空间通用的阈值。两种证明都把逐点的阈值变成有限个或一个收敛子列能控制的情形。

紧性是充分条件，不是必要条件：恒等映射在$\mathbb R$上均匀连续，而$\mathbb R$不紧。定理也只给出统一的阈值，不给出Lipschitz常数。$[0,1]$紧致，$\sqrt x$在其上连续因而均匀连续，但上一小节已说明它不Lipschitz连续。均匀连续保证存在随$\varepsilon$变化的$\delta(\varepsilon)$，Lipschitz连续则要求$\delta$可以取成$\varepsilon$的固定倍数。
